\subsection{CONNECTED SPACES AND CONNECTED SETS}

DEFINITION 1. A topological space X is said to be connected if it is not the union of two disjoint non-empty open sets.

An equivalent definition is obtained by replacing the words ``open sets'' by ``closed sets''. X is connected if and only if the only subsets of X which are both open and closed are the empty set and the whole space X.

If X is connected and if A, B are two non-empty open (resp. closed) subsets such that \(A \cup B = X\) , then \(A \cap B \neq \varnothing\) .

Examples. * 1) We shall see in Chapter IV, § 2, no. 5 that the real line is connected, and that the rational line is not. *

\begin{enumerate}
\def\labelenumi{\arabic{enumi})}
\setcounter{enumi}{1}
\tightlist
\item
  A discrete space which has more than one point is not connected.
\end{enumerate}

Observe that if \((\mathbf{U}_{t})_{t\in\mathbf{I}}\) is a partition of a topological space X consisting of open sets {[}necessarily non-empty, from the definition of a partition (Set Theory, R, § 4, no. 4){]} then each of the \(U_{t}\) is both open and closed in X, for the complement of \(U_{t}\) is the union of the \(U_{x}\) for which \(x\neq t\) . The open subsets of X are the subsets A such that \(A\cap U_{t}\) is open in \(U_{t}\) for each \(t\in I\) , so that X can be identified with the sum of the \(U_{t}\) (§ 2, no. 4, Example 3) and is not connected if I has more than one element.

DEFINITION 2. A subset A of a topological space X is said to be a connected set if the subspace A of X is connected.

For A to be a connected subset of X it is necessary and sufficient that, for each covering of A by two open (or closed) subsets B, C of X such that \(A \cap B\) and \(A \cap C\) are non-empty, we have \(A \cap B \cap C \neq \emptyset\) .

Examples. In any topological space, the empty set and every set consisting of a single point are connected. In a Hausdorff space X, every finite set consisting of more than one point is not connected, and more generally every subset of X which has more than one point and which has at least one isolated point is not connected.

If a dense subset A is connected, then the whole space X is connected; otherwise there would exist two non-empty disjoint open subsets M, N of X such that \(M \cup N = X\) , and \(M \cap A\) , \(N \cap A\) would be two disjoint non-empty open subsets of A whose union is A. Hence we have

PROPOSITION 1. If A is a connected set, then every set B such that \(A \subset B \subset \overline{A}\) is connected.

PROPOSITION 2. The union of a family of connected sets whose intersection is non-empty is connected.

Let \((\mathbf{A}_t)_{i\in I}\) be a family of connected subsets of \(\mathbf{X}\) , all of which contain the same point \(x\) ; we have to show that

\[
\mathbf {A} = \bigcup_ {i} \mathbf {A} _ {i}
\]

is connected. If not, there are two open sets B and C such that \(B \cap A\) and \(C \cap A\) are non-empty, and \(A \subset B \cup C\) and \(A \cap B \cap C = \varnothing\) . x belongs to one of the sets B, C, say \(x \in B\) ; on the other hand one of the sets \(A_t\) , say \(A_x\) , meets C; we have therefore \(A_x \subset B \cup C\) , \(A_x \cap B \cap C = \varnothing\) and \(B \cap A_x\) and \(C \cap A_x\) are non-empty. Hence \(A_x\) is not connected, which is a contradiction.

COROLLARY. Let \((\mathbf{A}_n)_{n\geqslant 0}\) be an infinite sequence of connected sets such that \(\mathbf{A}_{n + 1}\cap \mathbf{A}_n\neq \emptyset\) for all \(n\geqslant 0\) . Then the union \(\bigcup_{n = 0}^{\infty}\mathbf{A}_n\) is connected.

By induction on \(n\) we see immediately that the set \(\mathbf{B}_n = \bigcup_{i=0}^n \mathbf{A}_i\) is connected for all \(n\) , by Proposition 2. The sets \(\mathbf{B}_n\) have a non-empty intersection; hence their union, equal to \(\bigcup_{n=0}^{\infty} \mathbf{A}_n\) , is connected by Proposition 2.

PROPOSITION 3. Let A be a subset of a topological space X. If B is a connected subset of X which meets both A and \(\mathbb{C}\) A, then B meets the frontier of A.

For otherwise the intersections of B with the interior and exterior of A would be two open subsets of B which form a partition of B, and B would not be connected.

COROLLARY. In a connected space X, every non-empty set other than X itself has at least one frontier point.

